\scanpage{228}
% archive sheet number: 227

\mspage{- 5 -} \hfill n\textsuperscript{o}370

On définit de manière analogue les mesures invariantes, relativement invariantes et quasi-invariantes par $G$ sur $X$. Si $\mu$ est relativement invariante, on définit son multiplicateur $\chi$ par les formules

\[
\delta(s)\mu=\chi(s)\mu \tag*{(7')}
\]

\[
\int_X f(xs)\mathrm{d}\mu(x)=\chi(s)^{-1}\int_X f(x)\mathrm{d}\mu(x) \tag*{(8')}
\]

\[
\mu(As)=\chi(s)\mu(A). \tag*{(9')}
\]

Si on considère le groupe opposé $G^o$ à $G$ comme opérant dans $X$ par $(x,s)\to xs$, $\mu$ est relativement invariante par $G^o$ de même multiplicateur $\chi$.

Soit enfin $G$ un groupe localement compact. Il opère sur lui-même par translations à gauche et à droite, suivant les formules $\gamma(s)x=sx$, $\delta(s)x=xs^{-1}$. Tout ce qui précède est applicable, et on a donc, sur $G$, les notions de mesures \uline{invariantes à gauche}, \uline{invariantes à droite}, \uline{relativement invariantes à gauche}, \uline{relativement invariantes à droite}, \uline{quasi-invariantes à gauche}, \uline{quasi-invariantes à droite} (cf., toutefois, les n\textsuperscript{os} 8 et 9).

L'application $x\to x^{-1}$ est un homéomorphisme de $G$ sur $G$. Pour toute fonction $f$ sur $G$, on définira la fonction $\check f$ sur $G$ par

\[
\check f(x)=f(x^{-1}). \tag*{(10)}
\]

Pour toute mesure $\mu$ sur $G$, on définira la mesure $\check\mu$ par

\[
\check\mu(f)=\mu(\check f)\qquad\text{pour }f\in\mathcal{K}(G). \tag*{(11)}
\]

Autrement dit

\[
\int_G f(x)\mathrm{d}\check\mu(x)=\int_G f(x^{-1})\mathrm{d}\mu(x) \tag*{(12)}
\]

Si $A$ est un ensemble $\check\mu$-intégrable, $A^{-1}$ est $\mu$-intégrable, et

\[
\check\mu(A)=\mu(A^{-1}). \tag*{(13)}
\]
